\chapter{Introduction}\label{ch:intro}

Private messaging is now part of everyday personal, professional, and civic communication.
Signal~\cite{signal-home} is a widely used end-to-end encrypted messenger. Its cryptographic design
has influenced other messaging systems~\cite{signal-protocol-impact}. The Signal Protocol, introduced
in \cref{sec:bg-signal}, provides its end-to-end encryption. Signal's practical use and the protocol's
wider influence make it an important setting for studying key authentication in asynchronous private
messaging.

For journalists, activists, abuse survivors, and others whose safety may depend on confidential
communication, encryption alone is not sufficient. A message must be encrypted under the intended
recipient's authentic key. A silent key substitution defeats that guarantee while leaving the
conversation looking secure. An attacker may then read or alter communications without giving either
user an immediate indication that the channel has been compromised. Detecting such substitution is
therefore particularly important where a loss of confidentiality could cause serious harm.

End-to-end encryption protects message content only when it is encrypted to the intended recipient's
authentic key. In asynchronous messengers such as Signal, clients obtain that key from
provider-operated infrastructure. The public-key directory is therefore a critical trust point. If the
infrastructure substitutes an attacker's key for the recipient's, the protocol may establish an
encrypted channel with the attacker. This is a key-substitution, or ``fake-key,'' attack. The
encryption layer cannot distinguish the substituted key from the intended recipient's authentic key.

The traditional defence is a manual \emph{safety-number} comparison in which the two users in a
conversation compare a displayed code over a separate, out-of-band channel. The comparison is sound
when completed, but users rarely complete it without guidance. In one study, only about 14\% of
unguided pair--application trials succeeded. Just two of twelve pairs completed the comparison in any
of the applications tested~\cite{vaziripour2017}. Even a successful comparison is a one-off check that
says nothing about later key changes.

A messaging service's public-key directory maps user identities to their public identity keys.
\emph{Key transparency} (KT) addresses this limitation by committing directory states to a verifiable,
append-only history. Clients can then check their keys against a consistent record. Signal uses KT in
its optional Automatic Key Verification feature~\cite{signal-akv}.

For a substitution committed within the directory, existing systems still rely on the account holder
monitoring their own entry. A lookup can confirm that a fetched key matches the log, but this does not
tell the communicating peer that the logged key is wrong. No automatic, in-band channel warns that
peer. A returning client that monitors every epoch also incurs a cost that grows linearly with the
number of missed epochs.

Key verification in an asynchronous messenger should therefore operate automatically, warn both
parties in band, and reject a missing or invalid proof instead of accepting the key silently. A
returning client should also be able to verify missed history without downloading one proof for every
missed epoch, while keeping proof sizes and local verification costs practical.

\section{Aim and Approach}\label{sec:intro-aims}

We propose DUET (\emph{DU}al-index, \emph{E}poch-windowed \emph{T}ransparency), which performs
key-substitution checks automatically and reports warnings in band. The checks do not require a
separate manual comparison or change the asymptotic complexity of an individual lookup proof. The
complete mechanism nevertheless performs additional label-bound lookups and writes. DUET combines the
dual-index warning scheme with windowed consistency verification and is integrated into a Signal
Desktop prototype (\cref{ch:impl}). The windowed consistency mechanism replaces one proof per missed
epoch with a single proof whose size grows logarithmically in the endpoint log sizes.

DUET does not prevent a malicious provider from denying service. Its narrower goal is to ensure that a
substituted key, an invalid or missing proof, or an inconsistent directory history causes verification
to fail rather than allowing the client to accept an attacker's key silently. \Cref{ch:threat}
formalizes this distinction between detectable interference and availability.

\section{Contributions}\label{sec:intro-contributions}

This thesis makes four principal contributions.

\paragraph{Conceptual.} The \emph{dual-index warning scheme} derives a separate lookup label for each
party and a shared warning label from the conversation's shared secret. The warning label identifies a
shared slot to which either party can write an authenticated alert. For an honestly established and
monitored conversation, a malicious server cannot predict and pre-position that slot before the label
is first used. Before the warning commits, a censored write is detected through a key mismatch
alongside a proof-verified empty slot. After the warning commits, the transition audit detects any
removal as an invalid directory transition. The resulting signal allows both parties to detect the
substitution automatically and in band (\cref{sec:design-dual}).

The thesis also implements the abstract append-only-log component of the Authenticated Continuous Key
Agreement (ACKA) framework~\cite{dowling2023acka} as the DUET directory and evaluates its proof sizes
and local cryptographic costs. Within this directory, the dual-index scheme makes key substitution by
a semi-trusted operator detectable. The ACKA framework is introduced in \cref{sec:bg-acka}.

\paragraph{Cryptographic.} The cryptographic contribution applies an RFC~9162~\cite{rfc9162}
distant-head consistency proof to offline monitoring. A returning client re-verifies the log with one
proof instead of retrieving one proof for every missed epoch. The proof size stays logarithmic in the
endpoint log sizes rather than linear in the number of epochs missed (\cref{sec:design-windowed},
\cref{lem:windowed-size}). A same-head consistency lemma (\cref{lem:dualindex-consistency}) establishes
that Alice's peer lookup and Bob's self-check cannot authenticate different values under one signed
head. \Cref{ch:design} establishes these properties, and \cref{ch:eval} evaluates their costs.

\paragraph{Systems.} The systems contribution includes a standalone Rust reference implementation
comprising the \texttt{duet} crate, an HTTP directory server, and a cross-vantage auditor. It also
includes a TypeScript verifier integrated into a fork of the Signal Desktop client (\cref{ch:impl}).
The shared protocol operations in the two implementations are checked for byte-for-byte agreement
using cross-language test vectors.

\paragraph{Empirical.} The benchmark evaluation, covering metrics M1--M17, yields four principal
findings. The implemented VRF-indexed lookup proof is at rough size parity with CONIKS. The detection
channel uses additional lookups of the same shape rather than enlarging that proof. Median full-client
verification remains below \SI{150}{\micro\second} through $10^7$ labels and grows far more slowly than
$N$. A month-long offline catch-up uses a few hundred raw path-hash bytes, compared with hundreds of
kilobytes under the per-epoch baseline. The STRIDE and MITRE ATT\&CK evaluation distinguishes the
attacks detected by the implemented controls from the residual risks that remain (\cref{ch:eval}).

\section{Thesis Roadmap}\label{sec:intro-roadmap}

\Cref{ch:background} surveys the Signal stack, transparency logs, verifiable key directories, and the
ACKA framework, and derives the requirements addressed by DUET. \Cref{ch:threat} defines the threat
model and security requirements. \Cref{ch:design} presents the design and instantiates ACKA.
\Cref{ch:impl} describes the two-layer implementation and its Signal Desktop integration.
\Cref{ch:eval} evaluates the system empirically and against the threat model.
\Cref{ch:discussion,ch:lsepi,ch:conclusion} cover the discussion and future work, the legal, social,
ethical, and professional dimensions, and the conclusion. The appendices contain the formal
definitions, algorithms, implementation and test references, benchmark tables, security analyses,
traceability material, and glossary.
